\appendix

\section{Study design}

\begin{table}[!htbp]
  \caption{The four experiments, with the unit at which each is analysed.
  Frame pairs join two consecutive epochs at one standpoint and heading;
  transitions average the two headings of the same step.}
  \label{tab:design}
  \centering
  \scriptsize
  \begin{tabular}{@{}llllll@{}}
    \toprule
    Experiment & Source & Scope & Rows analysed & Unit of inference & Model \\
    \midrule
    Field panel  & Street View     & 5 US cities, 2007--2023 & 2{,}367 + 224 transitions & transition, clustered by site & \texttt{gemini-2.5-flash} \\
    Noise floor  & Streetscapes    & 150 images, 6 continents & 320 paired scorings & image & \texttt{gemini-2.5-flash} \\
    Perturbation & Streetscapes    & 150 images, 16 families & 9{,}750 responses & base image & three, see \S\ref{sec:probe} \\
    Optics       & Mapillary       & 40 panoramas, 3 conditions & 120 scored pairs & panorama & \texttt{gemini-2.5-flash} \\
    \bottomrule
  \end{tabular}
\end{table}

\section{Perturbation families and levels}

Sixteen families are applied at 57 levels in total: atmospheric scattering
($\beta = 0.5, 1.0, 2.0$), motion blur (5, 11, 21\,px), defocus (2, 4, 8\,px),
camera yaw (5, 10, 20$^\circ$), lens contamination (0.3, 0.6, 1.0), rain (0.3,
0.6, 1.0), glare (0.15, 0.35, 0.6), windscreen reflection (0.2, 0.4, 0.7),
foliage state (0.33, 0.66, 1.0), field of view (45, 60, 75$^\circ$), resolution
(0.5, 1.0, 2.1, 5.1, 8.3, 12.0\,MP), JPEG quality (25, 40, 60, 80, 95),
exposure ($\pm0.75$, $\pm1.5$\,EV), white balance (4000, 5000, 8000, 10000\,K),
contrast ($\times0.6, 0.8, 1.25, 1.6$) and saturation ($\times0.4, 0.7, 1.3,
1.8$). Levels span the range the corpus exhibits, with the atmospheric families
set by the calibration of Figure~\ref{fig:null}b. Camera yaw is applied as the
homography $K R K^{-1}$ for a pinhole source and by resampling the sphere for an
equirectangular one, so the manipulation is exact in both cases rather than
approximated by an image shift.

Rain was revised after calibration. It originally drew falling streaks; measured
against real rainy imagery, which is softer than clear imagery rather than
sharper, that mechanism moved the diagnostic statistic in the wrong direction,
and it was rewritten from the physical consequences of rain: overcast
illumination, a partly specular wet road, scattering in air, and droplets on the
lens.

\section{Prompt texts}

Five prompts are used. Two score perception, differing only in whether the
artefact block below is included. Three detect change: a task description alone,
the same with the artefact block, and the same again with the magnitude
threshold below. The two blocks are the experimental manipulation behind
Figure~\ref{fig:ladder} and Table~\ref{tab:threshold} and are given in full;
the surrounding scaffolding, which does not vary, is released with the code.

\paragraph{The artefact block.}
\begin{quote}\small\ttfamily
IMPORTANT --- judge the place, not the photograph.

These images come from crowdsourced street-level photography. They were taken
with different cameras, in different weather, at different times of day and
year, from slightly different positions. None of that is a property of the
street. Score the street as a competent observer standing there would perceive
it, discounting every one of the following:

--- image resolution, sharpness, compression artefacts, sensor noise \\
--- exposure, brightness, contrast, colour cast, white balance \\
--- motion blur, lens blur, lens flare, glare, water or dirt on the lens \\
--- weather at the moment of capture (rain, snow cover, fog, cloud) \\
--- time of day, sun angle, shadows, whether the scene is lit or overcast \\
--- season and the leaf state of the same plants \\
--- camera height, lens field of view, projection type, framing \\
--- vehicles, pedestrians, cyclists and other movable objects that happen to be
present

A blurry photograph of a well-maintained street is still a well-maintained
street. A crisp photograph of a derelict street is still a derelict street.
\end{quote}

\paragraph{The magnitude threshold.}
\begin{quote}\small\ttfamily
REPORT ONLY SUBSTANTIAL CHANGE.

A change is substantial if it alters what the street is or how it works, not
merely how it looks on the day. Report it when a structure, surface or fixed
element has been built, demolished, replaced, widened, narrowed or rebuilt.

Do not report: \\
--- graffiti, tags, posters, stickers or repainting of an existing surface \\
--- kerb or road markings repainted in the same configuration \\
--- construction hoarding, scaffolding, cones or temporary site barriers,
unless the building behind them has visibly changed \\
--- a plant that has grown, been trimmed, or changed with the season, where the
planting itself is unchanged \\
--- litter, dumped items, wear, staining, or a surface that is merely dirtier \\
--- individual items of street furniture appearing or disappearing

The test to apply: would a planner recording what happened on this street write
this down. If the honest answer is that the street is the same street with a
different surface appearance, report no change.
\end{quote}

\section{Released tables}

Figure~\ref{fig:ladder} averages over twelve dimensions in panel (a) and reports
them individually in panel (b). The underlying per-dimension figures, the signed
drift by dimension and by capture era, the decomposition of the pairwise
difference into between-standpoint and within-standpoint variance, the full
condition profiles behind Figure~\ref{fig:null}b, the family-by-level response
behind Figure~\ref{fig:surface}, and the city-by-city contrasts behind
Figure~\ref{fig:signal} are released as tables with the code.
